\begin{prob}
    Consider this version of quicksort given below. It is essentially the same
    as that given in lecture, except that 1) it always uses the last element of the
    array as the pivot, and 2) it has a \python{print} statement inserted at
    a crucial place.

    \inputminted{python}{\thisdir/include/quicksort.py}

    Suppose \python{arr} is an array of length $n$ with entries
    \python{[1, 2, 3, ..., n]}, where $n$ is some large integer.  If
    \python{quicksort(arr, 0, n)} is run, exactly how many times
    will
    \python{"hello"} be printed to the screen? Your answer should be an expression
    involving $n$, and should not involve $\sum$ or $\ldots$. Show your work.

    \begin{soln}
        Suppose we make a call to \python{quicksort} on a sorted array of size $k$.
        The pivot is set to the last (and largest) element in the array, therefore
        the condition of the \python{if}-statement in \python{partition} always
        evaluates to true, printing \python{"hello"} $k - 1$ times. 

        When the array is sorted, the root call to \python{quicksort} spawns
        two recursive calls; one on an array of size $n - 1$ and the other on
        an array of size zero. The first recursive call spawns two calls; one
        on an array of size $n - 2$ and the other on an array of size zero. And
        so forth. Only the recursive calls to arrays of size $> 1$ result in a call
        to \python{partition}. Therefore \python{partition} is called on arrays
        of size $n-1$, $n-2$, $n-3$, $\ldots$, $3$, $2$, and the total number
        of printed \python{"hello"}s is:
        \[
            (n-2) + (n-3) + \ldots + 3 + 2 + 1 = \frac{(n-1)n}{2}.
        \]
    \end{soln}

\end{prob}
